\documentclass[a4paper,10pt,fleqn]{article}
\usepackage[margin=1in]{geometry}
\usepackage{amssymb}
\usepackage{adjustbox}
\usepackage{natbib}
\usepackage[colorlinks=true,linkcolor=blue,citecolor=blue,urlcolor=blue]{hyperref}
\usepackage{fuzz}
\usepackage{schemapk}
\usepackage{zed-maths}
\usepackage{zed-proof}
\newdimen\savedleftskip
\title{Relational Calculus --- Music Library}
\author{txt2tex example}
\hypersetup{pdftitle={Relational Calculus --- Music Library},pdfauthor={txt2tex example},pdfcreator={txt2tex v1.6.0}}
\begin{document}

\maketitle

\begin{zed}[TrackId, ArtistId, AlbumId, PlaylistId, GenreName, TrackTitle, AlbumTitle, ArtistName]\end{zed}

\begin{schema}{Track}
tid : TrackId \\
title : TrackTitle \\
alid : AlbumId \\
duration : \nat \\
genre : GenreName
\end{schema}

\begin{schema}{Artist}
arid : ArtistId \\
name : ArtistName
\end{schema}

\begin{schema}{Album}
alid : AlbumId \\
title : AlbumTitle \\
arid : ArtistId
\end{schema}

\begin{axdef}
member : PlaylistId \rel TrackId
\end{axdef}

\section*{Query 1: Single-relation selection}

\noindent Retrieve the title and genre of every jazz track. This is a plain single-relation query over `Track` --- no join required.

\bigskip

\noindent
$\begin{array}{l}\{~ t : Track | t.genre = `Jazz' @ \\ \lblot~title == t.title, genre == t.genre~\rblot  ~\}\end{array}$


\section*{Query 2: Two-relation join via comprehension}

\noindent For each track, return its title together with the name of the artist who recorded the album it belongs to. Joins `Track` and `Album` on `AlbumId`, then `Album` and `Artist` on `ArtistId`.

\bigskip

\noindent
$\displaystyle
\begin{array}{l}
\begin{array}{l}\{~ t : Track ; al : Album ; ar : Artist | t.alid = al.alid \land \\
\quad al.arid = ar.arid @ \\ \lblot~track == t.title, artist == ar.name~\rblot  ~\}\end{array}
\end{array}$

\bigskip

\section*{Query 3: Three-relation chain with playlist}

\noindent List every track title that appears on at least one playlist, together with the album title. Chains `member`, `Track`, and `Album` through three declarations.

\bigskip

\noindent
$\displaystyle
\begin{array}{l}
\begin{array}{l}\{~ p : PlaylistId ; t : Track ; al : Album | (p, t.tid) \in member \land \\
\quad t.alid = al.alid @ \\ \lblot~playlist == p, track == t.title, album == al.title~\rblot  ~\}\end{array}
\end{array}$

\bigskip

\section*{Query 4: Binding output, filtered by duration}

\noindent Return a binding of `(track title, album title)` for every track whose duration exceeds 300 seconds. The characteristic expression is a two-component binding; no plain variable appears in the output.

\bigskip

\noindent
$\displaystyle
\begin{array}{l}
\begin{array}{l}\{~ t : Track ; al : Album | t.alid = al.alid \land \\
\quad t.duration > 300 @ \\ \lblot~track == t.title, album == al.title~\rblot  ~\}\end{array}
\end{array}$

\bigskip

\section*{Query 5: Multi-decl forall}

\noindent Assert that every track on any playlist belongs to a known album. Universal quantifier over two declarations with a conjunction predicate over projections.

\bigskip

\begin{zed}
\forall p : PlaylistId; t : Track | \\
\t1 (p, t.tid) \in member @ t.alid \in \{~ al : Album @ al.alid ~\}
\end{zed}

\section*{Query 6: Multi-decl comprehension, no characteristic expression}

\noindent Collect the set of `(Track, Album)` schema pairs for long tracks. When the characteristic expression is omitted Z defaults to the signature tuple (Z RM §3.9). The Q2(d) parser fix makes the dot-separator optional.

\bigskip

\noindent
$\displaystyle
\begin{array}{l}
\{~ t : Track ; al : Album | t.alid = al.alid \land \\
\quad t.duration > 300 ~\}
\end{array}$

\bigskip

\end{document}